\subsection{Materiales para el estudio de los datos}\label{sec:materiales_estudio}
\subsubsection{Hardware}

Para ejecutar los distintos entrenamientos se ha utilizado un ordenador con las siguientes características:

\begin{itemize}
    \item \textbf{CPU}: Intel Core i9-14900HX
    \item \textbf{RAM}: 32 GiB
    \item \textbf{GPU}: NVIDIA GeForce RTX 4060
    \item \textbf{GPU RAM}: 8 GiB
\end{itemize}

Estas características han resultado suficientes para completar los experimentos planteados sin limitaciones críticas.
El tiempo medio de entrenamiento por iteración ha sido aproximadamente 40 minutos, dependiendo de la configuración del modelo y de los parámetros de entrenamiento.

El tiempo total de cómputo se estima mediante:


\begin{equation*}
    \text{Tiempo total (h)} = \frac{N_{\text{iteraciones}} \times t_{\text{iteración}}}{60}
\end{equation*}

Donde \(N_{\text{iteraciones}} = 1917\) y \(t_{\text{iteración}} = 40 \text{minutos}\), resultando en un tiempo total aproximado de 766.8 horas (\(\approx 31 \text{días}\)).

El uso de hardware con mayores prestaciones permitiría reducir significativamente el tiempo total de entrenamiento y facilitar la realización de un mayor número de experimentos.

\subsubsection{Sofware}\label{sec:software}

Para el desarrollo del proyecto se ha utilizado una infraestructura basada en herramientas ampliamente empleadas en el ámbito del aprendizaje automático y la ingeniería de software.

Como lenguaje de programación principal se ha utilizado Python\footnote{https://www.python.org/}, debido a su amplio ecosistema de librerías científicas y de aprendizaje automático, así como su facilidad de integración con frameworks de alto nivel.

Para la construcción y entrenamiento de los modelos se han utilizado TensorFlow\footnote{https://www.tensorflow.org/}\par y Keras\footnote{https://keras.io/}, frameworks de aprendizaje profundo que permiten definir, entrenar y evaluar redes neuronales de forma eficiente. Keras se ha empleado como interfaz de alto nivel sobre TensorFlow, facilitando la experimentación rápida con distintas arquitecturas.

Durante las etapas iniciales del proyecto se empleó Jupyter Notebook\footnote{https://jupyter.org/} para la realización de pruebas exploratorias y la visualización preliminar de resultados. Sin embargo, su uso fue descartado en fases posteriores debido a la dificultad para mantener una estructura de código modular y escalable en experimentos de mayor complejidad.

Posteriormente, se adoptó Visual Studio Code\footnote{https://code.visualstudio.com/} como entorno de desarrollo principal, debido a su capacidad para gestionar múltiples tipos de archivos y su integración con herramientas de desarrollo. Se utilizaron extensiones específicas para Python, Jupyter y GitHub, permitiendo un flujo de trabajo unificado.

Como sistema de control de versiones se ha utilizado GitHub\footnote{https://github.com/}, que facilita la gestión del código, el seguimiento de cambios y la colaboración con el tutor del proyecto. Además, permite la sincronización del repositorio entre distintos dispositivos de forma eficiente. El código final desarrollado durante este trabajo se encuentra disponible en un repositorio público, referenciado en la bibliografía \cite{eeg_DanielRosiqueEgeaEeg}.

Para la visualización de resultados se ha utilizado Vega-Altair\footnote{https://altair-viz.github.io/}, una librería de visualización estadística declarativa enfocada en la simplicidad, eficiencia y legibilidad. A diferencia de enfoques imperativos como Matplotlib, Altair permite definir visualizaciones de forma estructurada y reproducible, enfocándose en lo que se quiere observar en lugar de en los detalles de la implementación. Esto conduce a un código más limpio y legible \cite{_IntroductionAltair2025}

Como gestor de base de datos se ha utilizado DuckDB\footnote{https://duckdb.org/}, un sistema de bases de datos analítico embebido que permite realizar consultas SQL de forma eficiente directamente sobre archivos locales. Su integración con Python y su alto rendimiento lo hacen especialmente adecuado para el procesamiento de grandes volúmenes de resultados experimentales \cite{_WhyDuckDB2026}.

Todas estas herramientas han sido integradas en un flujo de trabajo unificado orientado a la ejecución reproducible de experimentos.
